% 何昊程 — 中文简历（ATS 友好）
\documentclass[11pt,a4paper]{article}

% ── 页面布局 ─────────────────────────────────────────────────
\usepackage[margin=1.5cm,top=2cm,bottom=2cm]{geometry}
\usepackage{parskip}
\usepackage{titlesec}
\usepackage{enumitem}
\usepackage[hidelinks]{hyperref}

% ── 中文字体（xeCJK）────────────────────────────────────────
\usepackage{xeCJK}
\setCJKmainfont{Noto Sans CJK SC}
\setCJKsansfont{Noto Sans CJK SC}
\setCJKmonofont{Noto Sans CJK SC}

% ── 英文字体 ────────────────────────────────────────────────
\usepackage{fontspec}
\setmainfont{Latin Modern Roman}

% ── 去掉页码 ────────────────────────────────────────────────
\pagestyle{empty}

% ── 章节标题格式 ────────────────────────────────────────────
\titleformat{\section}{\large\bfseries}{}{0em}{}[\titlerule]
\titlespacing*{\section}{0pt}{12pt}{6pt}

% ── 列表格式 ────────────────────────────────────────────────
\setlist[itemize]{leftmargin=1.5em, itemsep=2pt, parsep=0pt, topsep=2pt}

% ── 自定义命令 ──────────────────────────────────────────────
\newcommand{\resumeItem}[1]{\item #1}
\newcommand{\experienceEntry}[4]{%
  \textbf{#1} \hfill #2 \\
  \textit{#3} \hfill \textit{#4}
}

\begin{document}

% ══════════════════════════════════════════════════════════════
%  联系方式
% ══════════════════════════════════════════════════════════════
\begin{center}
  {\LARGE\bfseries 何昊程} \\[6pt]
  haocheng.he@foxmail.com \enspace$|$\enspace
  中国香港特别行政区 \enspace$|$\enspace
  \href{https://linkedin.com/in/haocheng-he-92657028b}{linkedin.com/in/haocheng-he-92657028b}
\end{center}

% ══════════════════════════════════════════════════════════════
%  个人简介
% ══════════════════════════════════════════════════════════════
\section{个人简介}
工程学硕士，具备扎实的运筹学、数据分析与仿真建模能力。通过四段行业实习，
积累了航空运营、供应链优化及科技风险方面的实战经验。拥有多项发明专利与
软件著作权，期望从事运筹分析及科技风险相关岗位。

% ══════════════════════════════════════════════════════════════
%  教育背景
% ══════════════════════════════════════════════════════════════
\section{教育背景}

\textbf{香港理工大学} \hfill 中国香港 \\
\textit{航空工程学硕士} \hfill 2024年9月 -- 至今 \\
GPA：3.95 / 4.30

\vspace{6pt}

\textbf{广东工业大学} \hfill 广州 \\
\textit{工业工程学士} \hfill 2020年9月 -- 2024年6月 \\
GPA：3.67 / 4.00（前 20\%）

\vspace{2pt}

\textit{法学辅修} \hfill 2021年9月 -- 2024年6月 \\
GPA：3.61 / 4.00

% ══════════════════════════════════════════════════════════════
%  实习经历
% ══════════════════════════════════════════════════════════════
\section{实习经历}

\experienceEntry{运筹优化实习生}{2023年6月 -- 2023年9月}{供应链数据分析公司}{广州}
\begin{itemize}
  \resumeItem{使用 OR-Tools 构建仓库拣货路径优化模型，覆盖 3 个仓库，平均行走距离降低 18\%。}
  \resumeItem{基于 Python（pandas + matplotlib）自动化周报 KPI 看板，手动报表时间从 4 小时缩短至 20 分钟。}
  \resumeItem{分析逾 15 万条物流记录，识别瓶颈环节并提出布局优化方案，吞吐量提升 12\%。}
\end{itemize}

\vspace{4pt}

\experienceEntry{航空数据实习生}{2023年1月 -- 2023年5月}{机场管理局合作项目}{中国香港}
\begin{itemize}
  \resumeItem{使用 AnyLogic 开发航站楼旅客流离散事件仿真模型，测试 5 种布局方案。}
  \resumeItem{通过 SQL 处理超过 200 万条航班运行记录，构建延误预测特征集（准确率 82\%）。}
  \resumeItem{与 4 人团队协作完成可行性报告，获项目方采纳。}
\end{itemize}

\vspace{4pt}

\experienceEntry{生产仿真实习生}{2022年7月 -- 2022年10月}{汽车零部件制造企业}{广州}
\begin{itemize}
  \resumeItem{使用 Plant Simulation 搭建 12 工位装配线模型，定位导致 23\% 空闲时间的瓶颈工位。}
  \resumeItem{提出优化排班方案，在不增加人员的前提下将设备利用率从 71\% 提升至 84\%。}
  \resumeItem{撰写 30 页技术报告，记录仿真方法论，交付客户持续改善团队。}
\end{itemize}

\vspace{4pt}

\experienceEntry{商业智能实习生}{2021年6月 -- 2021年9月}{科技咨询公司}{深圳}
\begin{itemize}
  \resumeItem{设计并搭建 MySQL 数据仓库，整合 4 套遗留系统，支持 15+ 并发分析师查询。}
  \resumeItem{构建 Tableau 交互看板，追踪营收、成本及毛利率等高管关注的核心指标。}
  \resumeItem{编写 Python ETL 脚本，数据管线运行时间较人工流程缩短 40\%。}
\end{itemize}

% ══════════════════════════════════════════════════════════════
%  技能特长
% ══════════════════════════════════════════════════════════════
\section{技能特长}

\textbf{编程语言：}Python、Java、SQL \\
\textbf{工具与平台：}AnyLogic、Plant Simulation、OR-Tools、MySQL、Tableau、Git \\
\textbf{语言能力：}英语（雅思 6.5）、普通话（母语）、粤语（日常交流）

% ══════════════════════════════════════════════════════════════
%  学术成果与知识产权
% ══════════════════════════════════════════════════════════════
\section{学术成果与知识产权}

\begin{itemize}
  \resumeItem{发明专利 3 项（已授权）——物流与排程优化算法。}
  \resumeItem{实用新型专利 1 项——仓库设备人因工程设计。}
  \resumeItem{软件著作权 2 项——制造产线分析仿真工具包。}
  \resumeItem{优秀本科毕业论文——《基于仿真的混流装配线优化研究》。}
\end{itemize}

\end{document}
